%-------------------------
% Resume Builder shell: "Modern"
%
% Sans-serif section headings in a navy accent, a letter-spaced name, and a
% thin coloured rule under each heading. Body text stays roman for readability.
%
% CONTRACT — every shell must define \sectionsep and the \resume* macros used by
% backend/resume_render.py (\resumeItem, \resumeSubheading, \resumeProjectHeading,
% \resumeSubHeadingListStart/End, \resumeItemListStart/End) and expose the three
% {{...}} placeholders. Only the body between \begin{document} and \end{document}
% is generated; do NOT hand-edit the placeholders.
%------------------------

\documentclass[letterpaper,10pt]{article}

\usepackage{latexsym}
\usepackage[empty]{fullpage}
\usepackage{titlesec}
\usepackage{marvosym}
\usepackage[usenames,dvipsnames]{color}
\usepackage{verbatim}
\usepackage{enumitem}
\usepackage[hidelinks]{hyperref}
\usepackage{fancyhdr}
\usepackage[english]{babel}
\usepackage{tabularx}
% fontawesome5 is deliberately NOT loaded — its OTF fonts are unresolvable under
% Tectonic's XeTeX and abort the compile. See resume_templates.py.
\usepackage{multicol}
\usepackage{setspace}
\setstretch{1.02}
\setlength{\multicolsep}{-3.0pt}
\setlength{\columnsep}{-1pt}
% pdfTeX-only primitives; guarded so XeTeX doesn't abort. See classic.tex.
\ifdefined\pdfgentounicode
  \input{glyphtounicode}
\fi

% Accent used by the heading text and rule. Defined with \definecolor (not a
% package option) so no extra font/package resolution is involved.
\definecolor{accent}{RGB}{31,58,95}

\pagestyle{fancy}
\fancyhf{}
\fancyfoot{}
\renewcommand{\headrulewidth}{0pt}
\renewcommand{\footrulewidth}{0pt}

\addtolength{\oddsidemargin}{-0.55in}
\addtolength{\evensidemargin}{-0.55in}
\addtolength{\textwidth}{1.1in}
\addtolength{\topmargin}{-0.65in}
\addtolength{\textheight}{1.4in}

\urlstyle{same}

\raggedbottom
\raggedright
\setlength{\tabcolsep}{0in}

% Heading: sans-serif, accent-coloured, over a 0.5pt accent rule. \titlerule
% takes the rule height as an optional arg; the colour comes from the group.
\titleformat{\section}{
  \sffamily\raggedright\large\bfseries\color{accent}
}{}{0em}{}[{\color{accent}\titlerule[0.5pt]}\vspace{-4pt}]
\titlespacing*{\section}{0pt}{0pt}{8pt}

% See classic.tex for why the gap above a section is owned by \sectionsep and
% why \addvspace (not \vspace) is required here.
\newlength{\sectiongap}\setlength{\sectiongap}{8pt}
\newcommand{\sectionsep}{\par\addvspace{\sectiongap}}

\ifdefined\pdfgentounicode \pdfgentounicode=1 \fi

%-------------------------
% Layout macros — same contract as classic. Entry titles pick up the accent so
% the eye can scan employers/schools quickly.
\newcommand{\resumeItem}[1]{
  \item\small{
    {#1 \vspace{-2pt}}
  }
}

% Either ROW is emitted only when it has content. An entry may legitimately have
% no title (#1) or no location (#4) — the renderer always passes all four fields,
% empty string included — and without these guards such an entry rendered a blank
% line, leaving the dates or location floating alone beside empty space.
%
% Both rows are guarded because which row goes blank depends on which field is
% missing: the role is on row 1 and the company on row 2, so a title-less entry
% empties the FIRST row while a location-less one empties the second.
\newcommand{\resumeSubheading}[4]{
  \vspace{-2pt}\item
    \begin{tabular*}{1.0\textwidth}[t]{l@{\extracolsep{\fill}}r}
      \ifx\relax#1\relax\ifx\relax#2\relax\else
      \textbf{\color{accent}#1} & \textbf{\small #2} \\\fi\else
      \textbf{\color{accent}#1} & \textbf{\small #2} \\\fi
      \ifx\relax#3\relax\ifx\relax#4\relax\else
      \textit{\small#3} & \textit{\small #4} \\\fi\else
      \textit{\small#3} & \textit{\small #4} \\\fi
    \end{tabular*}\vspace{-7pt}
}

\newcommand{\resumeProjectHeading}[2]{
    \item
    \begin{tabular*}{1.001\textwidth}{l@{\extracolsep{\fill}}r}
      \small#1 & \textbf{\small #2}\\
    \end{tabular*}\vspace{-7pt}
}

% Education entries render through their own macro so a template CAN style them
% differently from experience (see compact.tex). This shell wants them identical,
% so it simply aliases the two.
\newcommand{\resumeEducationSubheading}[4]{\resumeSubheading{#1}{#2}{#3}{#4}}

\renewcommand\labelitemi{$\vcenter{\hbox{\tiny$\bullet$}}$}
\renewcommand\labelitemii{$\vcenter{\hbox{\tiny$\bullet$}}$}

\newcommand{\resumeSubHeadingListStart}{\begin{itemize}[leftmargin=0.0in, label={}]}
\newcommand{\resumeSubHeadingListEnd}{\end{itemize}}
\newcommand{\resumeItemListStart}{\begin{itemize}}
\newcommand{\resumeItemListEnd}{\end{itemize}\vspace{-5pt}}

%-------------------------------------------
\begin{document}

%----------HEADING----------
\begin{center}
    {\Huge \sffamily\bfseries \color{accent} {{FULL_NAME}}} \\ \vspace{4pt}
    \small {{CONTACT_LINE}}
    \vspace{-8pt}
\end{center}
% Offsets the header's trailing -8pt so the gap above the FIRST section matches
% the \sectionsep gap every later section gets. See classic.tex.
\par\vspace{1pt}

{{SECTIONS}}

\end{document}
